\begin{definition}
  Seja $M^{n}$ ($M$ superfície de dimenção $n$) uma superficie diferenciável.\\
  Dizemos que uma $n$-forma $Vol\in\Lambda^{n}(M)$ é uma \textbf{forma volume} para $M$ se para todo $p\in M$ tem-se que $Vol_p\neq 0$. 
\end{definition}
Localmente, uma forma volume sempre existe. Se $f:U\to V$ é uma parametrização, então
\begin{align*}
  Vol_V=dx^{1}\wedge\cdots\wedge dx^{n},
\end{align*}
é uma forma de volúme em $V\subseteq M$ e ela é chamada a \textbf{forma de volume local} em $M$ associada a parametrização $f$.
\begin{example}
  Se $B=\{e_1,\cdots,e_n\}$ é a base canônica de $\mathbb{R}^{n}$, então
  \begin{align*}
    Vol_{\mathbb{R}^{n}}=\det=e^{1}\wedge\cdots\wedge e^{n},
  \end{align*}
  é a forma volume para $\mathbb{R}^{n}$.
\end{example}
\begin{proposition}
  Se $F:M\to N$ é um difeomorfismo local é $Vol_{N}$ é a forma volume em $M$, então
  \begin{align*}
    Vol_{M}=F^{\ast}Vol_{N},
  \end{align*}
  é uma forma de volume em $M$.
\end{proposition}
\begin{proof} 
  \begin{align*}
    Vol_{M}(v)=F^{\ast}Vol_{N}(v)=(Vol_N)_{F(p)}(dF_{p}(v)),
  \end{align*}
  e como $dF_{p}$ é um isomorfismo, se $v=\{v_1,\cdots,v_n\}$ é uma base, então $dF_{p}(v)$ é uma base e como $Vol_N$ não se anula, temos que $Vol_M$ não se anula. 
\end{proof}
\begin{corollary}
  Se $\varphi:U\to V$ é uma parametrização, então
  \begin{align*}
    Vol_U=(\varphi^{-1})^{\ast}Vol_{\mathbb{R}^{n}}.
  \end{align*}
  Consequentemente, se
  \begin{align*}
    \varphi_1&:U_1\to V_1,\\
    \varphi_2&:U_2\to V_2,
  \end{align*}
  são outras parametrizações, temos que
  \begin{align*}
    Vol_{V_2}=det(\varphi^{-1}_2\circ\varphi_1)Vol_{V_1}.
  \end{align*}
\end{corollary}
\begin{proposition}
  Seja $M$ uma superfície diferenciável que possui uma forma volume $Vol$.\\
  Toda $n$-forma $\eta\in\Lambda^{n}(M)$ se escreve na forma
  \begin{align*}
    \eta=f Vol,
  \end{align*}
  para alguma função $f\in C^{\infty}(M)$.\\
  $\eta$ é também uma forma volume se, somente se $f$ não se anula em nenhum ponto $p\in M$. 
\end{proposition}
\begin{theorem}
  Uma superfície diferenciável $M$ é orientável se, somente se, ela admite uma forma volume.\\
  Além disso, esta forma determina uma orientação para $M$.
\end{theorem}
\begin{proof} 
  Exercício :c.  
\end{proof}
\begin{corollary}
  Se $M,N$ são superfícies diferenciáveis difeomorfas e $N$ é orientavel, então $M$ é orientável. 
\end{corollary}
\begin{definition}[Forma volume consistente]
   Seja $M$ uma superfície orientada e seja $\omega$ uma forma volume em $M$. Dizemos que $\omega$ é consistente com a orientação de $M$ se a orientação induzida por $\omega$ coincide com a orientação previamente fixada em $M$.
\end{definition}
